\begin{figure}[htbp]
\centering
\begin{tikzpicture}[
  box/.style={draw, rounded corners=2pt, align=center,
    text width=3.9cm, minimum height=3.2cm, inner sep=5pt,
    font=\small},
  flow/.style={-{Stealth[length=2mm]}, thick},
  note/.style={font=\footnotesize, align=center}
]
\node[box] (ambient) at (0,0) {
  \textbf{Independent arguments}\\[4pt]
  $a\in U$\\[3pt]
  $\Phi(a),\quad D_a\Phi$
};
\node[box] (graph) at (5.9,0) {
  \textbf{Proposed endpoints}\\[-1pt]
  \textbf{with argument ties}\\[4pt]
  $e=(e_-,e_+),\quad a=A(e)$\\[3pt]
  $\mathcal M(e)=\Phi(A(e))-z_+(e)$\\[3pt]
  {\footnotesize Nearby endpoints may mismatch}
};
\node[box] (matched) at (11.8,0) {
  \textbf{Matching endpoints}\\[4pt]
  $a=A(e)$\\[3pt]
  $\mathcal M(e)=0$
};
\draw[flow] (ambient.east) -- node[note, above=3pt]
  {impose\\$a=A(e)$} (graph.west);
\draw[flow] (graph.east) -- node[note, above=3pt]
  {impose\\$\mathcal M=0$} (matched.west);
\node[note] at (5.9,-2.05) {
  $D_e\mathcal M=(D_a\Phi)(A(e))\,D_eA-D_ez_+$
};
\end{tikzpicture}
\caption{One link near an actual passage, with physical parameters fixed.
The arrows impose successive equations.  Regular argument ties express
all independent kernel arguments, including coefficient copies, in terms
of nearby proposed endpoints before the scalar transfer equation is
imposed.  Here $z_+(e)$ is the prescribed output coordinate in the chosen
traversal.  On this local tied graph, $\mathcal M$ is the physical mismatch
in regular transverse coordinates, including at mismatched pairs.
Ambient partial derivatives are taken in $a$; endpoint derivatives use
the displayed chain rule.  This is a local statement, not a claim that
every point of the extended matching domain is a physical passage.}
\label{fig:matching-arguments}
\end{figure}
